\chapter{考察}
\label{chp:fifth}
\section{先行研究との違い}
\subsection{近似画像}
先行研究では，変化前が元画像である分，変化後の近似画像の再現度はかなり落ちてしまっているが，その一方で本研究は，元画像（先行研究の変化前）には及ばないものの十分な再現度を実現できている．
変化の前後での再現度のギャップもほとんど埋められており，本研究の目的の通りにパズルを生成することができた．
先行研究では実験されていなかったが，$n=20$であっても元の画像として十分に認識できるほど近似されていることも確認できた．

また，本研究では空所となる位置を，パリティ判定(\ref{fig:thm1}項)を用いてランダムに決定するのだが，
前章の実験で，そのような空所となる位置が存在せず，パズルが完成しないということは一つもなかったため，本研究においてこの決定法が非常に有効であったことがわかった．
\subsection{パズルの手数}
実験結果の図\ref{fig:n20},\ref{fig:n30},\ref{fig:n50}を見ると，先行研究の再現の方が少なかったり，本研究の方が少なかったりと，それらの数値にはそれほど差がなかったが，
全通りの組み合わせの手数の平均値を計算(\ref{tab:avepzl})し比較したところ，
$n=20,30,50$において，本研究の方が少なくなっていたため，\ref{sec:ippan}項の「距離コストを用いた行（列）の選択」によって，多少は改善されていたと考えることができる．
また，先行研究（再現）は，本研究のプログラムから「距離コストを用いた行（列）の選択」だけを除いた場合の手数なので，
先行研究より先行研究（再現）の方が手数が少ないことから，\ref{sec:ippan}項の「一般性をさらに持たせた移動手順の改善」によって，手数が少なくなったと考えられる．
\section{入力画像での違い}
図\ref{fig:ex8},\ref{fig:ex9}では，シマウマの模様の影響を受けてはいたものの，写真の濃淡や陰影などがうまく表現できていた．
図\ref{fig:ex3},\ref{fig:ex10}では，先行研究では白と黒のみの表現になってしまい，すこし見づらくなっていたが，合成駒を用いることでそれも改善できていた．
\section{今後の目的}
今後は，入力画像を2枚だけでなく3枚以上も可能にしたパズルを生成したい．
また，画像の近似値をさらに上げるために，グレースケール変換合成駒の生成に必要な計算式の改善もしていきたい．
